\documentclass[12pt,a4paper]{article}
\usepackage[utf8]{inputenc}
\usepackage[french]{babel}
\usepackage[T1]{fontenc}
\usepackage{geometry}
\usepackage{booktabs}
\usepackage{xcolor}
\usepackage{listings}
\usepackage{hyperref}
\usepackage{graphicx}
\usepackage{fancyhdr}
\usepackage{titlesec}
\usepackage{tcolorbox}

\geometry{margin=2.5cm}

% --- Colors and Styles ---
\definecolor{codegreen}{rgb}{0,0.6,0}
\definecolor{codegray}{rgb}{0.5,0.5,0.5}
\definecolor{codepurple}{rgb}{0.58,0,0.82}
\definecolor{backcolour}{rgb}{0.95,0.95,0.92}
\definecolor{primary}{RGB}{41, 128, 185}

\lstdefinestyle{mystyle}{
    backgroundcolor=\color{backcolour},   
    commentstyle=\color{codegreen},
    keywordstyle=\color{magenta},
    numberstyle=\tiny\color{codegray},
    stringstyle=\color{codepurple},
    basicstyle=\ttfamily\footnotesize,
    breakatwhitespace=false,         
    breaklines=true,                 
    captionpos=b,                    
    keepspaces=true,                 
    numbers=left,                    
    numbersep=5pt,                  
    showspaces=false,                
    showstringspaces=false,
    showtabs=false,                  
    tabsize=2
}
\lstset{style=mystyle}

\hypersetup{
    colorlinks=true,
    linkcolor=primary,
    filecolor=magenta,      
    urlcolor=blue,
}

\pagestyle{fancy}
\fancyhf{}
\rhead{\texttt{Disposition\_CTi}}
\lhead{Rapport Technique Avancé : Pipeline NLP}
\cfoot{\thepage}

\title{
    \vspace{-2cm}
    \textbf{\huge Architecture et Implémentation du Pipeline NLP}\\
    \large Analyse Technique Profonde avec Extraits de Code\\
    \vspace{0.5cm}
    \rule{\linewidth}{0.5mm}
}
\author{\textbf{Disposition\_CTi Development Team}}
\date{\today}

\begin{document}

\maketitle

\begin{abstract}
Ce rapport fournit une documentation technique complète de la phase d'enrichissement NLP (Natural Language Processing) de la plateforme \texttt{Disposition\_CTi}. Nous détaillons ici les mécanismes de bootstrap, l'analyse syntaxique pour la détection de rôles, et l'attribution sémantique locale, appuyés par des fragments du code source réel.
\end{abstract}

\tableofcontents
\newpage

\section{Phase 1 : Initialisation et Bootstrap NLTK}
Le pipeline commence par s'assurer que l'environnement est prêt. Contrairement à une simple recherche de texte, le pipeline NLP nécessite des modèles statistiques pré-entraînés.

\subsection{Gestion des Dépendances NLTK}
Le code suivant montre comment le système télécharge silencieusement les ressources nécessaires (\texttt{punkt}, \texttt{wordnet}, etc.) et bascule en mode "Basic" (Regex uniquement) si la connexion Internet est absente.

\begin{lstlisting}[language=Python, caption=Initialisation et Bootstrap]
def _bootstrap_nltk(self):
    self.nltk_enabled = True
    resources = ['punkt', 'averaged_perceptron_tagger', 
                 'maxent_ne_chunker', 'words', 'wordnet']
    for res in resources:
        try:
            nltk.download(res, quiet=True, raise_on_error=True)
        except Exception:
            self.nltk_enabled = False
            break
\end{lstlisting}

\section{Phase 2 : Analyse de Contexte et Détection de Rôle}
L'un des défis majeurs de la CTI est de savoir si une adresse IP mentionnée est l'attaquant ou une victime innocente. 

\subsection{Heuristique de Grammaire (POS Tagging)}
Le pipeline utilise le \textbf{Part-of-Speech Tagging} pour analyser les verbes entourant l'IOC. Si un IOC est suivi d'un verbe d'action comme \textit{"attacks"} ou \textit{"scans"}, il est identifié comme attaquant.

\begin{lstlisting}[language=Python, caption=Algorithme de Détection de Rôle]
# Extrait de nlp_enricher.py
for i in range(len(snippet)-1):
    word, tag = snippet[i]
    next_word, next_tag = snippet[i+1]
    # Si le mot courant est l'IOC et le suivant est un Verbe (VB)
    if entity.lower() in word.lower() and next_tag.startswith('VB'):
        if next_word.lower() in ["attacks", "scans", "targets", "drops"]:
            return {"role": "attacker", "confidence": 0.7}
\end{lstlisting}

\section{Phase 3 : Attribution Sémantique Locale}
Pour éviter de "polluer" tous les IOCs d'un rapport avec les mêmes tags, le système applique une attribution par phrase.

\subsection{Logique par Phrase}
Le texte est découpé en phrases. Les découvertes (familles de malwares, pays) faites dans la phrase "X" ne sont attribuées qu'aux IOCs présents dans cette même phrase "X".

\begin{lstlisting}[language=Python, caption=Attribution Sémantique Localisée]
def _attribute_findings(self, text, iocs):
    sentences = sent_tokenize(text)
    for sent in sentences:
        # 1. Identifier les IOCs presents dans CETTE phrase
        local_iocs = [ioc for ioc in iocs if ioc["value"] in sent]
        
        # 2. Extraire les categories uniquement de CETTE phrase
        categories = self._re_categories.findall(sent.lower())
        
        # 3. Attribuer les categories aux IOCs locaux uniquement
        for ioc in local_iocs:
            ioc["ioc_enrichment"]["threat_categories"] += categories
\end{lstlisting}

\section{Phase 4 : Résumé Extractif Avancé}
Pour les rapports longs (AlienVault, VirusTotal), le pipeline génère un résumé en calculant un score pour chaque phrase basé sur la fréquence des mots-clés techniques.

\section{Conclusion technique}
La structure modulaire (\texttt{BaseEnricher} $\rightarrow$ \texttt{NLPEnricher}) permet une maintenance aisée. L'utilisation hybride de Regex pour la rapidité et NLTK pour l'intelligence contextuelle offre un équilibre optimal pour un système de CTI moderne.

\end{document}
